\section{结论}
\label{sec:conclusion}

告警分诊要求本地模型完成两项工作：读取探针结果，以及据此做决定；而小模型只擅长前者。\hesp{} 据此分工：模型负责读，控制器选择探针、决定何时停止，并只接受账本证据支持的结论。在三个候选原因清单给定的模拟场景上，不含任何模型的控制器解决的案例与我们测试的每个模型一样多；而日志格式一变，固定解析器就完全失效，7B 读取器却没有。读取也正是攻击者写入的地方，所以 \hesp{} 永远不让模型自己的读数支持良性结论；有了这条规则、结束守卫和跨来源佐证，没有任何注入指令、伪造来源或针对读取器的自适应攻击能结掉一个案件。当候选原因无法事先得知，或证据只有组合起来才成立时，模型是否也应参与决定，需要固定程序解决不了的场景，以及我们考察过的公开安全数据集所不具备的数据。
